% Класификация по първия октет, начертана като дърво на решението.
%
% Рисува се отгоре надолу с явни координати, а не с автоматично разположение.
% Причината е, че трите клона имат различна дълбочина: TLS клонът стига до
% \code{ALPN}, клонът с открит текст стига до сравнение с преамбюла, а третият клон
% е лист още на първото ниво. Явните координати държат четирите листа на един ред.
%
% Черно на бяло: разликата се носи от формата (правоъгълник, ромб), от типа на
% линията (плътна, прекъсната) и от сивите запълвания. Никакъв цвят.
\begin{figure}[htbp]
\centering
\begin{tikzpicture}[
    scale=0.92, transform shape,
    % Листата са запълнени, за да се четат като край на път.
    leaf/.style={proc, fill=gray!25, minimum height=0.85cm},
    % Отхвърлената връзка е с прекъсната рамка и по-светло запълване: тя е
    % единственият изход, който не води до протокол.
    drop/.style={proc, fill=gray!10, densely dashed, thick, text width=2.6cm},
    lbl/.style={font=\fontsize{8}{9.5}\selectfont, align=left, inner sep=1.5pt,
                fill=white},
  ]

  % ---- корен -------------------------------------------------------------
  \node[box, text width=5.4cm, fill=gray!10] (root) at (6.8,7.8)
    {Първи октет от приетата връзка\\\code{read\_prefix}, поне един байт};

  % ---- ниво 1: трите изхода на \code{classify} ---------------------------
  \node[proc] (tls)   at (2.2,5.4) {TLS ръкостискане};
  \node[proc] (clear) at (8.5,5.4) {Открит текст};
  \node[drop] (drop)  at (13.0,5.4) {\code{Unknown}\\Връзката се затваря};

  % Общ хоризонтален сноп: и трите клона тръгват от една и съща точка, за да се
  % вижда, че това е едно решение с три изхода, а не три отделни проверки.
  \draw[flow] (root.south) -- ++(0,-0.55) -| (tls.north);
  \draw[flow] (root.south) -- ++(0,-0.55) -| (clear.north);
  \draw[flow] (root.south) -- ++(0,-0.55) -| (drop.north);

  \node[lbl, anchor=north west] at (2.32,6.62) {\code{0x16}\\TLS \code{ContentType}};
  \node[lbl, anchor=north west] at (8.62,6.62) {\code{[0x41, 0x5A]}\\главна буква ASCII};
  \node[lbl, anchor=south east] at (12.88,6.80) {друга стойност};

  \node[note, anchor=north, text width=3.0cm] at (13.0,4.75)
    {Не се гадае по следващите байтове.};

  % ---- ниво 2 ------------------------------------------------------------
  \node[proc] (alpn) at (2.2,4.0) {\code{ALPN} избира};
  \node[decision] (dec) at (8.5,4.0) {преамбюл на\\\code{HTTP/2}?};

  \draw[flow] (tls) -- (alpn);
  \draw[flow] (clear) -- (dec);

  % Сравнението е инкрементално. При съвпадащ префикс решението се отлага, вместо
  % да се чакат и 24-те байта на преамбюла. Оттам примката обратно към ромба.
  \draw[thin flow] (dec.north west) to[out=135, in=180, looseness=1.6] (dec.west);
  \node[note, anchor=south east] at (7.1,4.60) {\code{Maybe}: още байтове};

  % ---- листа -------------------------------------------------------------
  \node[leaf] (h2tls) at (0.6,1.9)  {\code{h2}\\над TLS};
  \node[leaf] (h1tls) at (3.5,1.9)  {\code{http/1.1}\\над TLS};
  \node[leaf] (h2pk)  at (6.9,1.9)  {\code{h2}\\открит текст};
  \node[leaf] (h11)   at (10.1,1.9) {\code{HTTP/1.1}\\открит текст};

  % Двата снопа са на една и съща височина, за да се четат четирите листа като
  % един ред, а не като два несвързани изхода.
  \coordinate (busl) at (2.2,2.90);
  \coordinate (busr) at (8.5,2.90);
  \draw[flow] (alpn.south) -- (busl) -| (h2tls.north);
  \draw[flow] (alpn.south) -- (busl) -| (h1tls.north);
  \draw[flow] (dec.south)  -- (busr) -| (h2pk.north);
  \draw[flow] (dec.south)  -- (busr) -| (h11.north);
  \node[lbl] at (7.55,2.90) {\code{Yes}};
  \node[lbl] at (9.55,2.90) {\code{No}};

  \node[note, anchor=north, text width=3.0cm] at (6.9,1.35)
    {Предварително знание, без \code{ALPN}.};
  \node[note, anchor=north, text width=3.0cm] at (10.15,1.35)
    {Надграждане до \code{WebSocket} или \code{h2c}.};

  % ---- твърдението, върху което стъпва решението с един байт ---------------
  \node[grp, note, text width=13.0cm, fill=gray!10] at (6.85,-0.35)
    {Двете начални множества не се пресичат. Точно един октет започва TLS запис,
     точно 26 октета започват метод в открит текст, нито един не прави и двете.
     Затова един байт стига. Проверено изчерпателно върху всичките 256 стойности
     на октета.};

\end{tikzpicture}
\caption{Класификация по първия октет на приетата връзка. Един прочетен байт
разделя TLS от открит текст, защото двете начални множества не се пресичат. При
открит текст инкрементално сравнение с преамбюла на \code{HTTP/2} разделя
\code{h2} с предварително знание от \code{HTTP/1.1}. Всяка друга стойност затваря
връзката, вместо да бъде отгатната.}
\label{fig:octet_classification}
\end{figure}
